\par\Needspace{7\baselineskip}\bigskip
\begin{center}
\pdfbookmark[2]{Gauss an Olbers, Göttingen, 15. Februar 1817}{pallas-brief-18170215}
\textsc{Gauss} an \textsc{Olbers}, Göttingen, 15. Februar 1817.
\end{center}
\noindent\hspace*{\parindent}. . . . .\quad Auch habe ich angefangen, für die Pallasstörungen eine Hülfstafel zu berechnen, eine Arbeit von ca. \(\frac{1}{2}\) Million Ziffern, und die ich ohne die thätige Beihülfe einiger jungen Leute, besonders des Hrn. \textsc{Westphal}, garnicht hätte unternehmen können. Mehr als die Hälfte ist schon fertig. Wenn die Tafel ganz vollendet ist, werde ich noch einmal alle Oppositionen auf das sorgfältigste berechnen und dabei auch die Beobachtungen von 1802, die bisher gar nicht mit angewandt waren, zuziehen. Die Anwendung der Wahrscheinlichkeitstheorie gibt \textls[180]{meiner} Bestimmung der {\AstroSymbols ♃}-Masse schon jetzt \textls[180]{dieselbe} Zuverlässigkeit, welche die \textsc{Bouvard}sche nach \textsc{Laplaces} Rechnung hat. Die letzte Opposition ist dabei noch nicht einmal zugezogen, sie wird den Werth nur sehr wenig ändern, aber die Zuverlässigkeit bedeutend vergrössern. Noch mehr wird die nächste {\AstroSymbols ☍} leisten. Suchen Sie doch auch die {\AstroSymbols ⚴} so bald wie möglich auf.\quad . . . . .
\par
